% compendium/portfolio-risk-financing/synthetic-shorts.tex
% Entry migrated on 2026-10-15 from the legacy securities-lending-and-synthetic-shorts file (slice 23), owned by
% synthetic-equity-cfd-trs (the file's map candidate); dispositions and source checks in notes/migration.md,
% the block verbatim in notes/artefacts/synthetic-equity-cfd-trs/derivations.tex.

\section{Synthetic short positions}

\begin{qfresult}[Synthetic short through options: long put plus short call is a short forward, not short stock, and does not remove the cost of scarce borrow]\label{res:synthetic-short-put-call}
\qffield{Statement} Take a European put and a European call on the same underlying, with the same strike $K$, maturity $T$ and contract multiplier, both settled the same way. Long the put and short the call pays at $T$
\[ \max(K-S_T,0)-\max(S_T-K,0)=K-S_T \]
for every terminal price $S_T$: a short forward position at delivery price $K$. With no arbitrage, a constant continuously compounded rate $r$ and a continuous dividend yield $q$, put--call parity gives the initial value
\[ P-C=Ke^{-rT}-S_0e^{-qT}=e^{-rT}(K-F),\qquad F=S_0e^{(r-q)T}, \]
so the package costs $e^{-rT}(K-F)$ upfront, positive or negative, and is a zero-cost short forward only when $K=F$.

Adding a \emph{short} zero-coupon bond that repays $K$ at $T$ (borrowing $Ke^{-rT}$ now) turns the package's payoff into $K-S_T-K=-S_T$. For a stock paying no dividends ($q=0$) that is the terminal value of one share sold short, and the initial values agree: $P-C-Ke^{-rT}=-S_0$. So long put, short call and short bond are equivalent to short stock; equivalently, long put and short call are equivalent to short stock plus a long bond paying $K$, a short sale whose proceeds are invested. With dividends the package less the bond is a terminal claim $-S_T$, worth $-S_0e^{-qT}$ now, while a share sold short is worth $-S_0$ and owes the dividends paid before $T$; the two differ by those dividend liabilities.

Harris's sidebar says that a long put, a short call with the same strike and maturity, ``plus a long position in a bond with the same maturity'', is economically equivalent to a short position in the underlying \citep[chap.~``Brokers'', sidebar ``Synthetic Short Positions'', p.~157, page image]{alpha-harris}. As worded, with a long bond paying $K$ the payoff is $2K-S_T$; the bond belongs on the other side, as above. The discrepancy is recorded here and Harris's wording is quoted as printed.

\qffield{Notation} $S_0$, $S_T$: underlying price now and at maturity; $P$, $C$: European put and call premia per unit of underlying (multiply by the contract multiplier for one contract); $F$: forward price; $r$: financing rate; $q$: dividend yield. With discrete dividends of present value $D$, replace $S_0e^{-qT}$ by $S_0-D$.

\qffield{Conditions} The identity needs matching strike, maturity, underlying and multiplier, the same settlement (both cash or both physical), European exercise, and frictionless borrowing and lending at $r$. Parity example (own computation, Black--Scholes prices with $S_0=100$, $K=95$, $r=5\%$, $q=2\%$, $\sigma=25\%$, $T=1$): $C=13.6847$, $P=6.0317$, $P-C=-7.6531=95e^{-0.05}-100e^{-0.02}$; $F=103.0455$ and $e^{-rT}(K-F)=-7.6531$. The script \nolinkurl{synthetic_short_checks.py} beside the artefact reprices both options by numerical integration over the risk-neutral density, and checks the payoffs leg by leg at $S_T=60, 95, 100, 140$: $K-S_T$ for the package, $-S_T$ with a short bond, $2K-S_T$ with a long bond (Interpretation).

\qffield{Practice} Against a short sale, Interpretation using the numbers above:
\begin{itemize}
\item \emph{package against short sale, no dividends}: take $q=0$, European options priced at frictionless parity, and borrowing or lending at $r$. Financing the premium $P-C=Ke^{-rT}-S_0$ at $r$, the package ends with $K-S_T-(P-C)e^{rT}=S_0e^{rT}-S_T$. Sell one share short, hold the whole proceeds $S_0$ as cash collateral credited at a constant rebate $r_b$ (the lending fee implicit in $r_b$, no other fee, no extra margin, no recall), and buy back at $T$: $S_0e^{r_bT}-S_T$. The two agree only when $r_b=r$, and the package is ahead by $S_0(e^{rT}-e^{r_bT})$ otherwise. With $S_0=100$, $K=95$, $r=5\%$, $\sigma=25\%$, $T=1$: $C=15.0471$, $P=5.4138$, $P-C=-9.6332$; at $S_T=100$ the package ends with $5.1271$ and the short sale with $5.1271$, $3.0455$ or $-0.9950$ for $r_b=5\%$, $3\%$ or $-1\%$, gaps of $0$, $2.0817$ and $6.1221$ at every $S_T$ (checked with day-by-day accounts in the script). The synthetic therefore replicates a short sale that earns the full market rebate, which is Harris's point: when options dealers receive the rebate and competition makes them price it in, the synthetic short ``in effect'' obtains short interest rebate (p.~157);
\item \emph{borrow cost}: a market maker who sells the put or buys the call is left long the stock's delta and hedges by selling the stock, short if it holds none (Interpretation), so a scarce borrow can be priced into the options and parity with a frictionless $r$ need not hold. That is an expectation, not a finding of the sources read; the package does not guarantee access to scarce borrow or remove its cost;
\item \emph{upfront and dividends}: the premium $P-C$ is paid or received at the start; dividends enter through $F$, whereas a short seller pays them to the lender as they fall due;
\item \emph{costs}: option spreads (\hyperref[meth:option-spread-measurement]{option spread measurement}) and the absence of listed options on many securities limit use; Harris notes that transaction costs ``can significantly reduce the economic benefits'' and that the strategy works only where options trade (p.~157).
\end{itemize}

\qffield{Pitfalls} Listed U.S. equity options are typically American and physically settled; that was not checked here against exchange or OCC documents. With American options the payoff identity holds only to maturity: the short call can be assigned early, notably before an ex-dividend date, which delivers a real short stock position that must then be borrowed or covered; the long put may be worth exercising early. Parity for American options is an inequality, not an equality. The legacy principle said the package ``replicate[s] a short forward'' and ``can replace scarce borrow, subject to funding, dividends, exercise, and option-liquidity costs''; the first is exact for European options, and the second is qualified above. Hull's treatment of put--call parity, a pathway reading, was not available and was not read; the derivation is standard and was checked numerically. Hu and Black note that a put and call with the same exercise price and expiration are economically equivalent to selling shares \citep[working-paper version, p.~5, footnote~6]{cfd-hu-black-2007}.

\qffield{See also} \hyperref[mech:securities-lending-short-sale]{securities lending for short sales}; \hyperref[def:securities-lending-fee-rebate]{lending fees and rebates}; \hyperref[res:option-market-maker-spreads]{option market-maker spreads}; \hyperref[def:arbitrage]{arbitrage}; \hyperref[res:arbitrage-risks]{arbitrage risks}.

\qffield{Further reading} The legacy block it replaces, verbatim, is in \nolinkurl{derivations.tex} in the folder \nolinkurl{notes/artefacts/synthetic-equity-cfd-trs/}.

\qffield{Index} synthetic short; put--call parity; long put short call; short forward; forward price; conversion and reversal; dividend yield; short interest rebate in options; American early exercise; assignment risk; ex-dividend early exercise; synthetic short with futures

\qffield{Status} Verification: partly checked 2026-10-15 (the text layer of \texttt{alpha-harris}, 2002 Oxford printing, p.~157; \texttt{cfd-hu-black-2007}, ECGI working-paper version, p.~5, footnote~6; the page image of p.~157 rechecked in a follow-up the same day; the identities derived and checked numerically by the script beside the artefact; \texttt{s7-hulloptions} not read). Stability: draft. Currency: timeless identity; market remarks as of 2002.
\end{qfresult}
